\chapter{String Algorithms}
\chapterepigraph{String algorithms are a small set of powerful machines --
the prefix and Z functions, Manacher's palindrome radii, polynomial
hashing, the suffix array, and the suffix automaton -- reused across many
problems. The CSES string section is really a tour of these five machines:
learn what each computes and the problems become a matter of feeding the
right string to the right machine.}

\section{The Five Machines}
\label{sec:cses-strings-intro}

Almost every problem here is solved by one of a handful of linear or
near-linear tools. Recognising which one a problem needs is the whole
game.

\begin{patternbox}[Which Machine Does the Problem Need?]
\begin{itemize}
  \item \textbf{Prefix function / Z-function} -- borders, periods, exact
        pattern matching, ``string built from repeated blocks'': Finding
        Borders, Finding Periods, String Matching, String Functions,
        Required Substring.
  \item \textbf{Manacher's algorithm} -- everything about palindromic
        substrings in $O(n)$: Longest Palindrome, All Palindromes.
  \item \textbf{Polynomial hashing} -- equality of substrings in $O(1)$
        after $O(n)$ preprocessing: Palindrome Queries, Repeating
        Substring, Finding Patterns.
  \item \textbf{Suffix array + LCP} -- ordering and counting substrings,
        reconstruction: Distinct Substrings, Inverse Suffix Array,
        Repeating Substring.
  \item \textbf{Suffix automaton} -- the Swiss-army knife for substring
        counting and ranking: Counting Patterns, Pattern Positions,
        Distinct Substrings, Substring Order I/II, Substring Distribution.
  \item \textbf{Standalone constructions} -- Minimal Rotation (Booth's),
        Word Combinations (trie + DP), Distinct Subsequences (DP), String
        Transform (inverse Burrows--Wheeler).
\end{itemize}
\end{patternbox}

\begin{ideabox}[Two Ideas Underlie Most of Them]
\textbf{Borders.} A border is a proper prefix that is also a suffix; the
prefix function stores the longest border of every prefix, and its
``failure chain'' enumerates all borders -- the basis of matching, periods,
and automaton building. \textbf{Endpos classes.} Two substrings are
equivalent if they occur at exactly the same set of ending positions; the
suffix automaton has one state per endpos class, which is why it counts
and ranks substrings so cheaply.
\end{ideabox}

Each section names the machine, explains why it applies, gives complete
compilable C++, a dry run, and complexity. Where a suffix automaton is
used, the same compact construction reappears so every listing stands
alone.

\newpage

\section{Word Combinations}
\label{sec:cses-word-combinations}

\problemstatement{Given a string $s$ of length $n$ and a dictionary of $k$
words, count the number of ways to build $s$ as a concatenation of
dictionary words (a word may be reused). Output the count modulo
$10^9+7$.}

\begin{patternbox}
\emph{``Count ways to segment a string into dictionary words''} is a
one-dimensional DP over prefix lengths, made fast by a \textbf{trie}. Let
$\textit{dp}[i]$ be the number of ways to form the first $i$ characters;
each way ends in some dictionary word.
\end{patternbox}

\subsection*{State and how the trie speeds it}

$\textit{dp}[i]$ = ways to form $s[0..i-1]$, with $\textit{dp}[0]=1$. From
a position $i$ with $\textit{dp}[i]>0$, walk down a trie of the dictionary
following $s[i],s[i+1],\ldots$; whenever the walk reaches a node marking
the end of a word (say the word spans $s[i..j-1]$), that word can extend
every way of forming the first $i$ characters, so
$\textit{dp}[j]\mathrel{+}=\textit{dp}[i]$. The trie lets us discover all
dictionary words starting at $i$ in time proportional to the longest such
word, without testing each dictionary word separately.

\begin{ideabox}[Push Contributions Forward Along the Trie]
Rather than asking ``which words end at $j$?'' we push: from each $i$ we
follow the single path $s[i..]$ in the trie and deposit $\textit{dp}[i]$
into $\textit{dp}[j]$ at every word-ending node. Each character of $s$ is
followed once per starting position, but the walk stops as soon as the
trie path dies.
\end{ideabox}

\subsection*{Algorithm}

\begin{enumerate}
  \item Insert all dictionary words into a trie, marking word ends.
  \item $\textit{dp}[0]=1$. For $i=0\ldots n-1$ with $\textit{dp}[i]>0$:
        walk the trie along $s[i..]$; at each word-end reached at column
        $j$, add $\textit{dp}[i]$ to $\textit{dp}[j]$.
  \item Answer $\textit{dp}[n]$.
\end{enumerate}

\subsection*{Dry run}

\dryrun{$s=$``ababc''; dictionary $\{$``ab'',``abab'',``c''$\}$.}

\begin{itemize}
  \item $\textit{dp}[0]=1$. From $0$: ``ab'' ends at $2\Rightarrow
        \textit{dp}[2]{+}{=}1$; ``abab'' ends at $4\Rightarrow
        \textit{dp}[4]{+}{=}1$.
  \item From $2$ ($\textit{dp}=1$): ``ab'' ends at $4\Rightarrow
        \textit{dp}[4]=2$.
  \item From $4$ ($\textit{dp}=2$): ``c'' ends at $5\Rightarrow
        \textit{dp}[5]=2$. Answer $2$.
\end{itemize}

\subsection*{C++ implementation}

\begin{lstlisting}
#include <bits/stdc++.h>
using namespace std;
const int MOD = 1e9 + 7;

struct Trie {
    struct Node { int nxt[26]; bool end; Node(){ memset(nxt,-1,sizeof nxt); end=false; } };
    vector<Node> t{Node()};
    void insert(const string& w) {
        int cur = 0;
        for (char c : w) {
            int x = c - 'a';
            if (t[cur].nxt[x] == -1) { t[cur].nxt[x] = t.size(); t.push_back(Node()); }
            cur = t[cur].nxt[x];
        }
        t[cur].end = true;
    }
};

int main() {
    string s; int k;
    cin >> s >> k;
    Trie tr;
    for (int i = 0; i < k; i++) { string w; cin >> w; tr.insert(w); }

    int n = s.size();
    vector<long long> dp(n + 1, 0);
    dp[0] = 1;
    for (int i = 0; i < n; i++) {
        if (!dp[i]) continue;
        int cur = 0;
        for (int j = i; j < n; j++) {
            int x = s[j] - 'a';
            if (tr.t[cur].nxt[x] == -1) break;      // no word continues
            cur = tr.t[cur].nxt[x];
            if (tr.t[cur].end) dp[j + 1] = (dp[j + 1] + dp[i]) % MOD;
        }
    }
    cout << dp[n] << "\n";
    return 0;
}
\end{lstlisting}

\begin{complexitybox}
\textbf{Time Complexity.} $O(n\cdot L)$ where $L$ is the longest word
(each start walks at most $L$ trie steps). \textbf{Space Complexity.}
$O(\text{total dictionary length}\cdot26)$ for the trie plus $O(n)$ DP.
\end{complexitybox}

\begin{takeawaybox}
Segmentation counting is prefix DP; a trie turns ``which words start
here?'' into a single character-by-character walk. The same push-forward
structure solves word-break existence, minimum-word segmentation, and
tokenisation.
\end{takeawaybox}

\section{String Matching}
\label{sec:cses-string-matching}

\problemstatement{Given a text $t$ and a pattern $p$, count how many times
$p$ occurs in $t$ (occurrences may overlap).}

\begin{patternbox}
\emph{``Count exact occurrences of a pattern''} is solved in linear time
by the \textbf{Z-function} (or KMP). Concatenate $p+\#+t$ with a separator
$\#$ not in either string; a position in the $t$-part whose Z-value equals
$|p|$ is an occurrence.
\end{patternbox}

\subsection*{The Z-function}

The Z-function of a string $w$ gives, for each index $i$, the length of the
longest substring starting at $i$ that matches a prefix of $w$. Build
$w=p+\#+t$. Wherever $Z[i]=|p|$ and $i$ lies in the $t$ portion, the text
matches the pattern's full length at that spot -- an occurrence. The
separator $\#$ guarantees a match cannot span the boundary. The Z-array
itself is computed in $O(|w|)$ using a maintained ``Z-box'' -- the
rightmost match interval seen so far.

\begin{ideabox}[Reuse Previous Matches via the Z-box]
When computing $Z[i]$, if $i$ falls inside a previously found match
interval $[l,r]$, the value $Z[i-l]$ is already known and can seed
$Z[i]$, so each character is compared a constant number of times
amortised. This is what makes Z (and KMP) linear rather than quadratic.
\end{ideabox}

\subsection*{Algorithm}

\begin{enumerate}
  \item Form $w=p+\#+t$ and compute its Z-array.
  \item Count indices $i$ (in the $t$ region) with $Z[i]=|p|$.
  \item Output the count.
\end{enumerate}

\subsection*{Dry run}

\dryrun{$t=$``saippua'', $p=$``pp''; $w=$``pp\#saippua''.}

\begin{itemize}
  \item $|p|=2$. Scanning $w$, the substring ``pp'' inside $t$ starts at
        one position; there $Z=2$.
  \item Exactly one index has $Z=2$, so the pattern occurs once.
\end{itemize}

\subsection*{C++ implementation}

\begin{lstlisting}
#include <bits/stdc++.h>
using namespace std;

vector<int> zFunction(const string& w) {
    int n = w.size();
    vector<int> z(n, 0);
    z[0] = n;
    int l = 0, r = 0;
    for (int i = 1; i < n; i++) {
        if (i < r) z[i] = min(r - i, z[i - l]);
        while (i + z[i] < n && w[z[i]] == w[i + z[i]]) z[i]++;
        if (i + z[i] > r) { l = i; r = i + z[i]; }
    }
    return z;
}

int main() {
    string t, p;
    cin >> t >> p;
    string w = p + "#" + t;
    vector<int> z = zFunction(w);

    int count = 0, m = p.size();
    for (int i = m + 1; i < (int)w.size(); i++)   // t region
        if (z[i] == m) count++;
    cout << count << "\n";
    return 0;
}
\end{lstlisting}

\begin{complexitybox}
\textbf{Time Complexity.} $O(|p|+|t|)$ -- one linear Z-array pass.
\textbf{Space Complexity.} $O(|p|+|t|)$ for the combined string and its
Z-array.
\end{complexitybox}

\begin{takeawaybox}
Exact pattern counting reduces to ``where does a prefix of $p+\#+t$
reappear at full pattern length,'' which the Z-function answers in linear
time. KMP's prefix function gives an equivalent solution; keep both in
your toolkit as interchangeable linear matchers.
\end{takeawaybox}

\section{Finding Borders}
\label{sec:cses-finding-borders}

\problemstatement{A border of a string is a proper prefix that is also a
suffix (shorter than the whole string). Given a string, output the lengths
of all its borders in increasing order.}

\begin{patternbox}
\emph{``All borders of a string''} is a direct application of the
\textbf{prefix function} (KMP failure function). The longest border of the
whole string is $\pi[n-1]$; following the ``failure chain''
$\pi[n-1]\to\pi[\pi[n-1]-1]\to\cdots$ enumerates every border.
\end{patternbox}

\subsection*{Prefix function and the failure chain}

The prefix function $\pi[i]$ is the length of the longest proper border of
the prefix $s[0..i]$. A key fact: the set of all border lengths of the
\emph{whole} string is exactly $\{\pi[n-1],\ \pi[\pi[n-1]-1],\ \ldots\}$
down to $0$. This is because the longest border's own longest border is
the next-longest border of the whole string, and so on. Collecting the
chain and reversing it yields all borders in increasing order.

\begin{ideabox}[Borders Nest Like Russian Dolls]
Every border of $s$ is itself bordered by the next-shorter border. So the
failure chain from $\pi[n-1]$ visits all of them in decreasing order --
each obtained by looking up the prefix function at the previous border's
last index.
\end{ideabox}

\subsection*{Algorithm}

\begin{enumerate}
  \item Compute the prefix function $\pi$ of $s$.
  \item Starting from $b=\pi[n-1]$, repeatedly record $b$ and set
        $b=\pi[b-1]$ until $b=0$.
  \item Reverse the collected lengths and output them (increasing order).
\end{enumerate}

\subsection*{Dry run}

\dryrun{$s=$``abcababcab''.}

\begin{itemize}
  \item The whole string's longest border is ``abcab'' (length $5$):
        $\pi[9]=5$.
  \item $\pi[4]=2$ (border ``ab''); $\pi[1]=0$. Chain: $5,2$.
  \item Increasing order: $2,5$ -- the border lengths.
\end{itemize}

\subsection*{C++ implementation}

\begin{lstlisting}
#include <bits/stdc++.h>
using namespace std;

vector<int> prefixFunction(const string& s) {
    int n = s.size();
    vector<int> pi(n, 0);
    for (int i = 1; i < n; i++) {
        int j = pi[i - 1];
        while (j > 0 && s[i] != s[j]) j = pi[j - 1];
        if (s[i] == s[j]) j++;
        pi[i] = j;
    }
    return pi;
}

int main() {
    string s;
    cin >> s;
    int n = s.size();
    vector<int> pi = prefixFunction(s);

    vector<int> borders;
    for (int b = pi[n - 1]; b > 0; b = pi[b - 1]) borders.push_back(b);
    reverse(borders.begin(), borders.end());

    for (int b : borders) cout << b << " ";
    cout << "\n";
    return 0;
}
\end{lstlisting}

\begin{complexitybox}
\textbf{Time Complexity.} $O(n)$ for the prefix function; the failure
chain has $O(n)$ total length. \textbf{Space Complexity.} $O(n)$.
\end{complexitybox}

\begin{takeawaybox}
Borders are the prefix function's reason for existing. Walking the failure
chain from $\pi[n-1]$ lists them all. This chain is also how KMP recovers
after a mismatch and how periods (next section) are derived.
\end{takeawaybox}

\section{Finding Periods}
\label{sec:cses-finding-periods}

\problemstatement{A period of a string $s$ of length $n$ is an integer
$p$ ($1\le p\le n$) such that $s[i]=s[i+p]$ for all valid $i$ (equivalently,
$s$ is a prefix of an infinite repetition of its first $p$ characters).
Output all periods in increasing order.}

\begin{patternbox}
Periods and borders are two sides of one coin: \textbf{$p$ is a period of
$s$ if and only if $s$ has a border of length $n-p$}. So compute the
prefix function, enumerate all border lengths $b$ (including $0$), and
each gives a period $n-b$.
\end{patternbox}

\subsection*{The period--border duality}

If $s$ has a border of length $b$, the first $n-b$ characters repeat to
cover the string, making $n-b$ a period; conversely every period $p$
leaves a border of length $n-p$. The border lengths of the whole string
are the failure chain $\pi[n-1],\pi[\pi[n-1]-1],\ldots$ together with $0$
(the empty border, which corresponds to the trivial period $p=n$). Mapping
each border length $b$ to $n-b$ and sorting gives all periods.

\begin{ideabox}[Every Border Length $b$ Gives Period $n-b$]
Include the empty border ($b=0$) so that $p=n$ (the whole string is always
its own period) is present. The failure chain supplies the remaining
borders. One prefix-function pass yields every period.
\end{ideabox}

\subsection*{Algorithm}

\begin{enumerate}
  \item Compute the prefix function $\pi$.
  \item Collect border lengths via the failure chain from $\pi[n-1]$,
        plus $0$.
  \item For each border length $b$, output $n-b$; sort ascending.
\end{enumerate}

\subsection*{Dry run}

\dryrun{$s=$``abcabca'' ($n=7$).}

\begin{itemize}
  \item Longest border ``abca'' has length $4$; chain gives borders
        $4,1$; plus $0$.
  \item Periods: $7-4=3$, $7-1=6$, $7-0=7$.
  \item Increasing: $3,6,7$ (indeed ``abc'' repeats with period $3$).
\end{itemize}

\subsection*{C++ implementation}

\begin{lstlisting}
#include <bits/stdc++.h>
using namespace std;

vector<int> prefixFunction(const string& s) {
    int n = s.size();
    vector<int> pi(n, 0);
    for (int i = 1; i < n; i++) {
        int j = pi[i - 1];
        while (j > 0 && s[i] != s[j]) j = pi[j - 1];
        if (s[i] == s[j]) j++;
        pi[i] = j;
    }
    return pi;
}

int main() {
    string s;
    cin >> s;
    int n = s.size();
    vector<int> pi = prefixFunction(s);

    vector<int> periods;
    for (int b = pi[n - 1]; b > 0; b = pi[b - 1]) periods.push_back(n - b);
    periods.push_back(n);                         // border 0 -> period n
    sort(periods.begin(), periods.end());

    for (int p : periods) cout << p << " ";
    cout << "\n";
    return 0;
}
\end{lstlisting}

\begin{complexitybox}
\textbf{Time Complexity.} $O(n)$. \textbf{Space Complexity.} $O(n)$.
\end{complexitybox}

\begin{takeawaybox}
Period equals $n$ minus a border length. Once you see that borders and
periods are complementary, every ``how does this string repeat?'' question
becomes a prefix-function query. This duality also underlies detecting the
smallest repeating block of a string.
\end{takeawaybox}

\section{Minimal Rotation}
\label{sec:cses-minimal-rotation}

\problemstatement{Given a string $s$, consider all $n$ rotations
(cyclic shifts) of it. Output the lexicographically smallest rotation.}

\begin{patternbox}
\emph{``Lexicographically smallest cyclic rotation''} has a linear-time
solution -- \textbf{Booth's algorithm} (or the equivalent least-rotation
scan) -- that finds the starting index of the minimal rotation by
comparing candidate starts and skipping ranges that cannot win.
\end{patternbox}

\subsection*{The least-rotation scan}

Work on $s+s$ (so every rotation is a length-$n$ window). Maintain two
candidate start positions $i$ and $j$ and an offset $k$: compare
$s[i+k]$ with $s[j+k]$. If they are equal, extend $k$. If one is smaller,
the losing candidate -- and, crucially, \emph{every start it skipped over}
-- can be discarded by jumping it past $j+k$ (or $i+k$). Because each
comparison either advances $k$ or eliminates a range of candidates, the
total work is linear. The surviving candidate is the minimal rotation's
start.

\begin{ideabox}[Skip Doomed Starts in Bulk]
When candidate $j$ beats candidate $i$ after matching $k$ characters,
none of the starts strictly between $i$ and $i+k+1$ can be minimal either
(they were dominated), so we jump $i$ past them all. This bulk elimination
is what keeps the scan $O(n)$ despite comparing rotations.
\end{ideabox}

\subsection*{Algorithm}

\begin{enumerate}
  \item Set $i=0$, $j=1$, $k=0$ over the doubled string $ss=s+s$.
  \item While all three are in range, compare $ss[i+k]$ and $ss[j+k]$:
        equal $\Rightarrow k{+}{+}$; otherwise advance the larger
        candidate past $\max(i,j)+k+1$, reset $k=0$, and keep $i<j$.
  \item The answer starts at $\min(i,j)$; output $ss[\text{start}\,..\,
        \text{start}+n-1]$.
\end{enumerate}

\subsection*{Dry run}

\dryrun{$s=$``acab''.}

\begin{itemize}
  \item Rotations: ``acab'',``caba'',``abac'',``baca''.
  \item Lexicographically smallest is ``abac'', starting at index $2$.
  \item The scan converges on start $=2$.
\end{itemize}

\subsection*{C++ implementation}

\begin{lstlisting}
#include <bits/stdc++.h>
using namespace std;

int leastRotation(const string& s) {
    string ss = s + s;
    int n = ss.size();
    int i = 0, j = 1, k = 0;
    while (i < (int)s.size() && j < (int)s.size() && k < (int)s.size()) {
        char a = ss[i + k], b = ss[j + k];
        if (a == b) { k++; }
        else if (a > b) { i = max(i + k + 1, j + 1); k = 0; }
        else            { j = max(j + k + 1, i + 1); k = 0; }
    }
    return min(i, j);                             // start of minimal rotation
}

int main() {
    string s;
    cin >> s;
    int start = leastRotation(s);
    string ss = s + s;
    cout << ss.substr(start, s.size()) << "\n";
    return 0;
}
\end{lstlisting}

\begin{complexitybox}
\textbf{Time Complexity.} $O(n)$ -- each comparison advances $k$ or a
candidate pointer, and both are bounded by $O(n)$ total. \textbf{Space
Complexity.} $O(n)$ for the doubled string.
\end{complexitybox}

\begin{takeawaybox}
The least-rotation scan finds the minimal cyclic shift in linear time by
eliminating dominated starts in bulk. It is the canonical tool for
canonicalising cyclic strings -- useful for necklace counting and
comparing circular sequences up to rotation.
\end{takeawaybox}

\section{Longest Palindrome}
\label{sec:cses-longest-palindrome}

\problemstatement{Given a string, find its longest substring that is a
palindrome (reads the same forwards and backwards). Output that
substring.}

\begin{patternbox}
\emph{``Longest palindromic substring''} is exactly what \textbf{Manacher's
algorithm} computes in $O(n)$: for every possible centre it finds the
radius of the longest palindrome centred there, reusing already-known
radii via a maintained ``current rightmost palindrome.''
\end{patternbox}

\subsection*{Manacher in one idea}

For each centre, a palindrome extends symmetrically until the characters
disagree. The naive approach re-expands every centre from scratch
($O(n^2)$). Manacher keeps the palindrome with the farthest right edge
seen so far, centred at $c$ with right boundary $r$. For a new centre $i<r$,
its mirror $2c-i$ already has a known radius, which can seed $i$'s radius
(bounded by $r-i$) before any character comparison -- so each character is
matched only a constant number of times amortised. Handling even-length
palindromes is done either with a second pass or by inserting separators.

\begin{ideabox}[Mirror Symmetry Seeds Each New Centre]
Inside the current rightmost palindrome, position $i$ mirrors $2c-i$, so
its radius starts at least as large as the mirror's (capped by the
boundary). Only growth beyond the boundary requires real comparisons,
which advance $r$ -- giving linear total work.
\end{ideabox}

\subsection*{Algorithm}

\begin{enumerate}
  \item Transform $s$ to $t$ by inserting a sentinel (e.g.\ \texttt{\#})
        between every character and at the ends, so all palindromes become
        odd-length.
  \item Compute radii $P[i]$ over $t$ using the mirror-seed rule and a
        maintained $(c,r)$.
  \item The maximum $P[i]$ locates the longest palindrome; map back to
        indices in $s$ and output the substring.
\end{enumerate}

\subsection*{Dry run}

\dryrun{$s=$``aybabtu''.}

\begin{itemize}
  \item Centred at ``bab'', the palindrome extends to radius covering
        ``bab''.
  \item No longer palindrome exists, so the answer is ``bab'' (length
        $3$).
\end{itemize}

\subsection*{C++ implementation}

\begin{lstlisting}
#include <bits/stdc++.h>
using namespace std;

int main() {
    string s;
    cin >> s;
    // transform: ^ # a # b # ... $   (sentinels avoid bounds checks)
    string t = "^";
    for (char c : s) { t += '#'; t += c; }
    t += "#$";

    int n = t.size();
    vector<int> P(n, 0);
    int c = 0, r = 0, bestLen = 0, bestCenter = 0;
    for (int i = 1; i < n - 1; i++) {
        if (i < r) P[i] = min(r - i, P[2 * c - i]);   // mirror seed
        while (t[i + P[i] + 1] == t[i - P[i] - 1]) P[i]++;
        if (i + P[i] > r) { c = i; r = i + P[i]; }
        if (P[i] > bestLen) { bestLen = P[i]; bestCenter = i; }
    }
    int start = (bestCenter - bestLen) / 2;         // map back to s
    cout << s.substr(start, bestLen) << "\n";
    return 0;
}
\end{lstlisting}

\begin{complexitybox}
\textbf{Time Complexity.} $O(n)$ -- each character is matched a constant
number of times amortised. \textbf{Space Complexity.} $O(n)$ for the
transformed string and radius array.
\end{complexitybox}

\begin{takeawaybox}
Manacher computes every palindromic radius in linear time by seeding new
centres from their mirrors within the current rightmost palindrome. The
sentinel transform unifies odd and even cases. It is the definitive tool
for palindrome-substring problems, reused in the next section to count
them all.
\end{takeawaybox}

\section{All Palindromes}
\label{sec:cses-all-palindromes}

\problemstatement{Given a string, count the total number of palindromic
substrings (each occurrence counted, for every start/end pair that forms a
palindrome).}

\begin{patternbox}
\emph{``Count all palindromic substrings''} is again \textbf{Manacher's
algorithm}: the radius at each centre directly gives how many palindromes
are centred there, so summing the radii counts every palindromic substring
in $O(n)$.
\end{patternbox}

\subsection*{Radii are counts}

For an odd centre at character $i$ with maximal radius $d_1[i]$ (counting
the centre itself), there are exactly $d_1[i]$ odd-length palindromes
centred at $i$: lengths $1,3,\ldots,2d_1[i]-1$. Likewise each even centre
between two characters contributes $d_2[i]$ even-length palindromes. So
the total number of palindromic substrings is
$\sum_i d_1[i]+\sum_i d_2[i]$. Manacher computes both radius arrays in
linear time, and the sum is the answer.

\begin{ideabox}[Every Nested Palindrome Counts Once]
A palindrome of radius $k$ around a centre contains, at the same centre,
palindromes of every smaller radius -- $k$ of them. Adding the radius at
each centre therefore counts every (centre, length) palindrome exactly
once, with no double counting across centres.
\end{ideabox}

\subsection*{Algorithm}

\begin{enumerate}
  \item Run Manacher on the sentinel-transformed string to get radii
        $P[i]$.
  \item In the transformed indexing, each centre contributes
        $\lceil P[i]/2\rceil$ palindromic substrings; sum over all
        centres (a single combined pass counts both odd and even).
  \item Output the total.
\end{enumerate}

\subsection*{Dry run}

\dryrun{$s=$``aaa''.}

\begin{itemize}
  \item Single characters: $3$ palindromes. Length-2 ``aa'': $2$.
        Length-3 ``aaa'': $1$.
  \item Total $=3+2+1=6$, which equals the sum of Manacher radii.
\end{itemize}

\subsection*{C++ implementation}

\begin{lstlisting}
#include <bits/stdc++.h>
using namespace std;

int main() {
    string s;
    cin >> s;
    string t = "^";
    for (char c : s) { t += '#'; t += c; }
    t += "#$";

    int n = t.size();
    vector<int> P(n, 0);
    int c = 0, r = 0;
    long long total = 0;
    for (int i = 1; i < n - 1; i++) {
        if (i < r) P[i] = min(r - i, P[2 * c - i]);
        while (t[i + P[i] + 1] == t[i - P[i] - 1]) P[i]++;
        if (i + P[i] > r) { c = i; r = i + P[i]; }
        total += (P[i] + 1) / 2;                  // palindromes at this center
    }
    cout << total << "\n";
    return 0;
}
\end{lstlisting}

\begin{complexitybox}
\textbf{Time Complexity.} $O(n)$ -- one Manacher pass. \textbf{Space
Complexity.} $O(n)$ for the transformed string and radii.
\end{complexitybox}

\begin{takeawaybox}
The radius Manacher computes at each centre \emph{is} the count of
palindromes there, so counting all palindromic substrings is a single sum.
For counting \emph{distinct} palindromes instead, an Eertree (palindromic
tree) is the specialised structure -- but for total occurrences, Manacher
suffices.
\end{takeawaybox}

\section{Required Substring}
\label{sec:cses-required-substring}

\problemstatement{Count the number of strings of length $n$ over the
lowercase alphabet (26 letters) that contain a given pattern $p$ as a
substring. Output the count modulo $10^9+7$.}

\begin{patternbox}
\emph{``Count length-$n$ strings that \textbf{contain} a pattern''} is
easier via the complement: count strings that \emph{avoid} $p$, then
subtract from $26^n$. Avoidance is tracked by a DP running the
\textbf{KMP automaton} -- the state is how much of $p$ is currently
matched.
\end{patternbox}

\subsection*{DP over KMP automaton states}

Let $m=|p|$. Build the KMP automaton: from a state $j$ (meaning the last
$j$ characters matched a prefix of $p$) reading character $c$, the
automaton moves to $\delta(j,c)$, the new longest matched prefix. Define
$\textit{dp}[i][j]$ = number of length-$i$ strings whose automaton state is
$j$ \emph{without ever reaching state $m$}. Transition over all 26
characters:
\[
  \textit{dp}[i+1][\delta(j,c)]\mathrel{+}=\textit{dp}[i][j],\quad
  \text{only when }\delta(j,c)<m.
\]
Reaching state $m$ means $p$ appeared, so those transitions are dropped
from the ``avoid'' count. The answer is
$26^n-\sum_{j<m}\textit{dp}[n][j]$.

\begin{ideabox}[Match Progress Is the Only State That Matters]
Whether a string avoids $p$ depends solely on how much of $p$ is currently
matched -- the KMP state. Building the automaton once ($O(26m)$) lets the
DP forget the actual characters and track just the $m$ possible progress
levels.
\end{ideabox}

\subsection*{Algorithm}

\begin{enumerate}
  \item Compute $p$'s prefix function and the transition $\delta(j,c)$
        for $j=0\ldots m-1$, $c\in\{a..z\}$.
  \item $\textit{dp}[0][0]=1$. For $i=0\ldots n-1$, push each state over
        all 26 characters, discarding transitions into state $m$.
  \item Answer $26^n-\sum_{j=0}^{m-1}\textit{dp}[n][j]$ modulo $10^9+7$.
\end{enumerate}

\subsection*{Dry run}

\dryrun{$n=2$, $p=$``aa''; alphabet size $2$ for illustration ($\{a,b\}$).}

\begin{itemize}
  \item Strings of length $2$: $aa,ab,ba,bb$. Only $aa$ contains ``aa''.
  \item Avoidance DP counts $3$ ($ab,ba,bb$); total $2^2-3=1$. With the
        real 26-letter alphabet the same logic gives the count.
\end{itemize}

\subsection*{C++ implementation}

\begin{lstlisting}
#include <bits/stdc++.h>
using namespace std;
const long long MOD = 1e9 + 7;

int main() {
    int n;
    string p;
    cin >> n >> p;
    int m = p.size();

    vector<int> pi(m, 0);
    for (int i = 1; i < m; i++) {
        int j = pi[i - 1];
        while (j > 0 && p[i] != p[j]) j = pi[j - 1];
        if (p[i] == p[j]) j++;
        pi[i] = j;
    }
    // automaton transitions delta[j][c], j in 0..m
    vector<array<int,26>> delta(m + 1);
    for (int j = 0; j <= m; j++)
        for (int c = 0; c < 26; c++) {
            if (j < m && c == p[j] - 'a') delta[j][c] = j + 1;
            else delta[j][c] = (j == 0) ? 0 : delta[pi[j - 1]][c];
        }

    vector<vector<long long>> dp(n + 1, vector<long long>(m + 1, 0));
    dp[0][0] = 1;
    for (int i = 0; i < n; i++)
        for (int j = 0; j < m; j++) if (dp[i][j])
            for (int c = 0; c < 26; c++) {
                int nj = delta[j][c];
                if (nj < m) dp[i + 1][nj] = (dp[i + 1][nj] + dp[i][j]) % MOD;
            }

    long long total = 1;                           // 26^n
    for (int i = 0; i < n; i++) total = total * 26 % MOD;
    long long avoid = 0;
    for (int j = 0; j < m; j++) avoid = (avoid + dp[n][j]) % MOD;
    cout << ((total - avoid) % MOD + MOD) % MOD << "\n";
    return 0;
}
\end{lstlisting}

\begin{complexitybox}
\textbf{Time Complexity.} $O(n\cdot m\cdot26)$ for the DP, plus
$O(m\cdot26)$ to build the automaton. \textbf{Space Complexity.}
$O(n\cdot m)$ (reducible to $O(m)$ with rolling rows).
\end{complexitybox}

\begin{takeawaybox}
``Contains a pattern'' is best counted as $26^n$ minus ``avoids the
pattern,'' and avoidance is a DP over KMP-automaton states. Turning a
pattern into an automaton and running a counting DP on its states is a
recurring competitive-programming template.
\end{takeawaybox}

\section{Palindrome Queries}
\label{sec:cses-palindrome-queries}

\problemstatement{Given a string, process two kinds of queries: (1) change
the character at position $k$; (2) ask whether the substring $[a,b]$ is a
palindrome. Answer all type-2 queries online.}

\begin{patternbox}
\emph{``Update a character, then test a substring for palindrome''}
combines point updates with range equality checks -- exactly what
\textbf{polynomial hashing over a Fenwick tree} supports. Maintain a
forward hash and a reversed hash; the substring is a palindrome iff the
two agree.
\end{patternbox}

\subsection*{Two hashes, both point-updatable}

Store the string's characters weighted by powers of a base $B$ in two
Fenwick (BIT) trees: one for the forward direction, one for the reversed
direction. The hash of $[a,b]$ read forwards is a normalised range sum
from the forward BIT; the hash of $[a,b]$ read backwards is the
corresponding range sum from the reversed BIT. Normalising both to the
same starting power lets them be compared directly. A character update is
two point updates (one per tree). The substring is a palindrome exactly
when the forward and reversed hashes match.

\begin{ideabox}[Palindrome = Forward Hash Equals Reversed Hash]
A substring reads the same both ways iff its forward and backward hashes
coincide. Encoding each character as $c\cdot B^{\text{pos}}$ in a BIT makes
both hashes range sums, so point updates and range queries are each
$O(\log n)$. Use two independent moduli/bases to make hash collisions
negligible.
\end{ideabox}

\subsection*{Algorithm}

\begin{enumerate}
  \item Precompute powers of $B$; insert each character into a forward
        and a reversed Fenwick tree.
  \item Update: point-update both trees at the changed position.
  \item Query $[a,b]$: fetch the forward and reversed range hashes,
        normalise to a common power, compare; equal $\Rightarrow$
        palindrome.
\end{enumerate}

\subsection*{Dry run}

\dryrun{$s=$``aybabtu''; query palindrome on ``bab'' ($[2,4]$).}

\begin{itemize}
  \item Forward hash of ``bab'' and reversed hash of ``bab'' both encode
        $b,a,b$ in mirror order and match.
  \item After updating position $3$ (`a'$\to$`x'), ``bxb'' still matches
        (palindrome); ``bxa'' would not.
\end{itemize}

\subsection*{C++ implementation}

\begin{lstlisting}
#include <bits/stdc++.h>
using namespace std;
typedef unsigned long long ull;
const ull B = 131;

struct BIT {
    int n; vector<ull> tree;
    BIT(int n) : n(n), tree(n + 1, 0) {}
    void update(int i, ull val) { for (++i; i <= n; i += i & -i) tree[i] += val; }
    ull query(int i) { ull s = 0; for (++i; i > 0; i -= i & -i) s += tree[i]; return s; }
    ull range(int l, int r) { return query(r) - (l ? query(l - 1) : 0); }
};

int main() {
    int n, q;
    string s;
    cin >> n >> q >> s;
    vector<ull> pw(n + 1, 1);
    for (int i = 1; i <= n; i++) pw[i] = pw[i - 1] * B;

    BIT fwd(n), rev(n);
    for (int i = 0; i < n; i++) {
        fwd.update(i, (ull)s[i] * pw[i]);          // weight by position
        rev.update(i, (ull)s[i] * pw[n - 1 - i]);
    }

    while (q--) {
        int type; cin >> type;
        if (type == 1) {
            int k; char c; cin >> k >> c; k--;
            fwd.update(k, ((ull)c - s[k]) * pw[k]);
            rev.update(k, ((ull)c - s[k]) * pw[n - 1 - k]);
            s[k] = c;
        } else {
            int a, b; cin >> a >> b; a--; b--;
            ull h1 = fwd.range(a, b) * pw[n - 1 - b];       // normalise
            ull h2 = rev.range(a, b) * pw[a];
            cout << (h1 == h2 ? "YES" : "NO") << "\n";
        }
    }
    return 0;
}
\end{lstlisting}

\begin{complexitybox}
\textbf{Time Complexity.} $O((n+q)\log n)$ -- each update and query is
$O(\log n)$. \textbf{Space Complexity.} $O(n)$ for the two Fenwick trees
and power table.
\end{complexitybox}

\begin{takeawaybox}
Dynamic substring equality is polynomial hashing plus a point-updatable
Fenwick tree; a palindrome test is just ``forward hash equals reversed
hash.'' For safety against adversarial collisions, run two independent
hashes. This forward/reversed-hash pattern answers many online substring
comparison queries.
\end{takeawaybox}

\section{Finding Patterns}
\label{sec:cses-finding-patterns}

\problemstatement{Given a text $t$ and $k$ patterns, determine for each
pattern whether it appears as a substring of $t$. Output ``YES'' or
``NO'' per pattern.}

\begin{patternbox}
\emph{``Do these many patterns occur in one fixed text?''} is the
signature use of a \textbf{suffix automaton (SAM)} built on $t$. The SAM
recognises exactly the substrings of $t$, so each pattern is a matter of
walking its characters through the automaton.
\end{patternbox}

\subsection*{The suffix automaton recognises every substring}

A suffix automaton of $t$ is the smallest automaton whose paths from the
initial state spell precisely the substrings of $t$. Building it takes
$O(|t|)$ (with a constant or $\log\Sigma$ factor for transitions). To test
a pattern $p$, start at the initial state and follow the transition for
each character of $p$; if at any point the needed transition is missing,
$p$ is not a substring, otherwise it is. The total query cost is the sum
of pattern lengths.

\begin{ideabox}[Build Once, Query Many]
The SAM front-loads all the work into a single linear construction over
$t$; afterwards every pattern is answered in time proportional to its own
length. This ``preprocess the text, stream the queries'' shape is why SAM
dominates multi-pattern substring questions.
\end{ideabox}

\subsection*{Algorithm}

\begin{enumerate}
  \item Build the suffix automaton of $t$ by extending it one character at
        a time.
  \item For each pattern, walk from the initial state following its
        characters; report YES if the walk completes, else NO.
\end{enumerate}

\subsection*{Dry run}

\dryrun{$t=$``aybabtu''; patterns ``bab'', ``ba'', ``xyz''.}

\begin{itemize}
  \item ``bab'': walk b$\to$a$\to$b succeeds -- YES.
  \item ``ba'': succeeds -- YES.
  \item ``xyz'': the transition on ``x'' is missing at the start -- NO.
\end{itemize}

\subsection*{C++ implementation}

\begin{lstlisting}
#include <bits/stdc++.h>
using namespace std;

struct SAM {
    struct State { int len, link; map<char,int> next; };
    vector<State> st;
    int last;
    SAM() { st.push_back({0, -1, {}}); last = 0; }
    void extend(char c) {
        int cur = st.size();
        st.push_back({st[last].len + 1, -1, {}});
        int p = last;
        while (p != -1 && !st[p].next.count(c)) { st[p].next[c] = cur; p = st[p].link; }
        if (p == -1) st[cur].link = 0;
        else {
            int q = st[p].next[c];
            if (st[p].len + 1 == st[q].len) st[cur].link = q;
            else {
                int clone = st.size();
                st.push_back(st[q]);
                st[clone].len = st[p].len + 1;
                while (p != -1 && st[p].next[c] == q) { st[p].next[c] = clone; p = st[p].link; }
                st[q].link = st[cur].link = clone;
            }
        }
        last = cur;
    }
};

int main() {
    string t;
    int k;
    cin >> t >> k;
    SAM sam;
    for (char c : t) sam.extend(c);

    while (k--) {
        string p;
        cin >> p;
        int cur = 0;
        bool ok = true;
        for (char c : p) {
            auto it = sam.st[cur].next.find(c);
            if (it == sam.st[cur].next.end()) { ok = false; break; }
            cur = it->second;
        }
        cout << (ok ? "YES" : "NO") << "\n";
    }
    return 0;
}
\end{lstlisting}

\begin{complexitybox}
\textbf{Time Complexity.} $O(|t|\log\Sigma)$ to build, plus
$O(\sum|p_i|\log\Sigma)$ for all queries. \textbf{Space Complexity.}
$O(|t|)$ states.
\end{complexitybox}

\begin{takeawaybox}
A suffix automaton turns ``is this a substring?'' into a walk, so many
patterns against a fixed text are answered by one construction plus cheap
traversals. The same SAM, augmented with counts, answers the next two
problems (Counting Patterns, Pattern Positions) with almost no extra code.
\end{takeawaybox}

\section{Counting Patterns}
\label{sec:cses-counting-patterns}

\problemstatement{Given a text $t$ and $k$ patterns, count for each
pattern how many times it occurs in $t$ (occurrences may overlap).}

\begin{patternbox}
This is Finding Patterns with counts: the number of occurrences of a
substring equals the size of its \textbf{endpos} set, which the suffix
automaton stores per state. Compute those sizes once, then each pattern's
count is a lookup after walking to its state.
\end{patternbox}

\subsection*{Endpos sizes count occurrences}

In a suffix automaton, each state represents a set of substrings that
share the same set of ending positions (their \emph{endpos} set). The
number of occurrences of any substring in a state equals that state's
endpos-set size. To compute these: mark each state created as a genuine
prefix-endpoint (the ``\textit{last}'' after each extension) with count
$1$ and clones with $0$, then propagate counts up the suffix-link tree in
order of decreasing length -- each state adds its count to its link. A
pattern's count is the endpos size of the state its walk ends in.

\begin{ideabox}[Propagate Counts Up the Link Tree]
The suffix-link tree groups substrings by endpos containment: a child's
occurrences are a subset that also count toward the parent. Summing counts
from longer states into their links (processing by decreasing $\textit{len}$)
computes every endpos size in one linear pass.
\end{ideabox}

\subsection*{Algorithm}

\begin{enumerate}
  \item Build the SAM of $t$; set $\textit{cnt}=1$ for each non-clone
        state (the \textit{last} after each extend), $0$ for clones.
  \item Sort states by $\textit{len}$ descending; add each state's
        $\textit{cnt}$ to its suffix link's $\textit{cnt}$.
  \item For each pattern, walk to its state; output that state's
        $\textit{cnt}$ (or $0$ if the walk fails).
\end{enumerate}

\subsection*{Dry run}

\dryrun{$t=$``aaa''; pattern ``a''.}

\begin{itemize}
  \item ``a'' occurs at positions $1,2,3$ -- endpos size $3$.
  \item The SAM state for ``a'' accumulates count $3$ after propagation,
        which is the reported answer.
\end{itemize}

\subsection*{C++ implementation}

\begin{lstlisting}
#include <bits/stdc++.h>
using namespace std;

struct SAM {
    struct State { int len, link; map<char,int> next; };
    vector<State> st; vector<long long> cnt; int last;
    SAM() { st.push_back({0, -1, {}}); cnt.push_back(0); last = 0; }
    void extend(char c) {
        int cur = st.size();
        st.push_back({st[last].len + 1, -1, {}}); cnt.push_back(1);  // real endpoint
        int p = last;
        while (p != -1 && !st[p].next.count(c)) { st[p].next[c] = cur; p = st[p].link; }
        if (p == -1) st[cur].link = 0;
        else {
            int q = st[p].next[c];
            if (st[p].len + 1 == st[q].len) st[cur].link = q;
            else {
                int clone = st.size();
                st.push_back(st[q]); cnt.push_back(0);              // clone: no endpoint
                st[clone].len = st[p].len + 1;
                while (p != -1 && st[p].next[c] == q) { st[p].next[c] = clone; p = st[p].link; }
                st[q].link = st[cur].link = clone;
            }
        }
        last = cur;
    }
};

int main() {
    string t; int k;
    cin >> t >> k;
    SAM sam;
    for (char c : t) sam.extend(c);

    int m = sam.st.size();
    vector<int> order(m);
    for (int i = 0; i < m; i++) order[i] = i;
    sort(order.begin(), order.end(),
         [&](int a, int b){ return sam.st[a].len > sam.st[b].len; });
    for (int v : order)
        if (sam.st[v].link >= 0) sam.cnt[sam.st[v].link] += sam.cnt[v];

    while (k--) {
        string p; cin >> p;
        int cur = 0; bool ok = true;
        for (char c : p) {
            auto it = sam.st[cur].next.find(c);
            if (it == sam.st[cur].next.end()) { ok = false; break; }
            cur = it->second;
        }
        cout << (ok ? sam.cnt[cur] : 0) << "\n";
    }
    return 0;
}
\end{lstlisting}

\begin{complexitybox}
\textbf{Time Complexity.} $O(|t|\log\Sigma)$ build and count propagation,
plus $O(\sum|p_i|\log\Sigma)$ queries. \textbf{Space Complexity.}
$O(|t|)$.
\end{complexitybox}

\begin{takeawaybox}
Occurrence count = endpos-set size, computed by propagating counts up the
suffix-link tree. This single augmentation upgrades the SAM from ``does it
occur?'' to ``how many times?'' -- and the same endpos machinery drives
distinct-substring counting and ranking later in the chapter.
\end{takeawaybox}

\section{Pattern Positions}
\label{sec:cses-pattern-positions}

\problemstatement{Given a text $t$ and $k$ patterns, output for each
pattern the starting position of its \emph{first} occurrence in $t$
(1-indexed), or $-1$ if it does not occur.}

\begin{patternbox}
Same suffix automaton, one more annotation: store, for each state, the
ending position of the \textbf{first} occurrence of its substrings
(\textit{firstpos}). A pattern's first-occurrence start is then
$\textit{firstpos}-\lvert p\rvert+1$.
\end{patternbox}

\subsection*{Tracking the first endpos}

When the automaton is extended with the character at index $i$ (1-indexed),
the newly created state's longest substring first \emph{ends} at position
$i$, so set its $\textit{firstpos}=i$. A cloned state inherits the
$\textit{firstpos}$ of the state it was split from, because the clone
represents shorter substrings that first appear no later. After building,
a pattern $p$ that walks to state $v$ has its earliest occurrence ending at
$\textit{firstpos}[v]$, hence starting at
$\textit{firstpos}[v]-\lvert p\rvert+1$.

\begin{ideabox}[The First Endpoint Locates the First Occurrence]
Every substring in a state first occurs at the same earliest ending
position (that is what an endpos class guarantees for its shortest common
context). Recording that position at creation time -- and copying it to
clones -- gives $O(1)$ first-occurrence lookups after the walk.
\end{ideabox}

\subsection*{Algorithm}

\begin{enumerate}
  \item Build the SAM; when creating the state at step $i$, set
        $\textit{firstpos}=i$; for a clone, copy the original's
        $\textit{firstpos}$.
  \item For each pattern, walk to its state $v$; if the walk fails, output
        $-1$; else output $\textit{firstpos}[v]-\lvert p\rvert+1$.
\end{enumerate}

\subsection*{Dry run}

\dryrun{$t=$``aybabtu''; pattern ``ab''.}

\begin{itemize}
  \item ``ab'' first occurs ending at position $4$ (characters $3$--$4$).
  \item Start $=4-2+1=3$, the 1-indexed first position of ``ab''.
\end{itemize}

\subsection*{C++ implementation}

\begin{lstlisting}
#include <bits/stdc++.h>
using namespace std;

struct SAM {
    struct State { int len, link, firstpos; map<char,int> next; };
    vector<State> st; int last;
    SAM() { st.push_back({0, -1, 0, {}}); last = 0; }
    void extend(char c, int pos) {               // pos = 1-indexed end position
        int cur = st.size();
        st.push_back({st[last].len + 1, -1, pos, {}});
        int p = last;
        while (p != -1 && !st[p].next.count(c)) { st[p].next[c] = cur; p = st[p].link; }
        if (p == -1) st[cur].link = 0;
        else {
            int q = st[p].next[c];
            if (st[p].len + 1 == st[q].len) st[cur].link = q;
            else {
                int clone = st.size();
                st.push_back(st[q]);             // inherits firstpos of q
                st[clone].len = st[p].len + 1;
                while (p != -1 && st[p].next[c] == q) { st[p].next[c] = clone; p = st[p].link; }
                st[q].link = st[cur].link = clone;
            }
        }
        last = cur;
    }
};

int main() {
    string t; int k;
    cin >> t >> k;
    SAM sam;
    for (int i = 0; i < (int)t.size(); i++) sam.extend(t[i], i + 1);

    while (k--) {
        string p; cin >> p;
        int cur = 0; bool ok = true;
        for (char c : p) {
            auto it = sam.st[cur].next.find(c);
            if (it == sam.st[cur].next.end()) { ok = false; break; }
            cur = it->second;
        }
        cout << (ok ? sam.st[cur].firstpos - (int)p.size() + 1 : -1) << "\n";
    }
    return 0;
}
\end{lstlisting}

\begin{complexitybox}
\textbf{Time Complexity.} $O(|t|\log\Sigma)$ build, $O(\sum|p_i|
\log\Sigma)$ queries. \textbf{Space Complexity.} $O(|t|)$.
\end{complexitybox}

\begin{takeawaybox}
Recording each state's first ending position turns the SAM into a
first-occurrence locator. Combined with endpos counts (previous section),
one automaton answers presence, frequency, and position -- the full
substring query trio.
\end{takeawaybox}

\section{Distinct Substrings}
\label{sec:cses-distinct-substrings}

\problemstatement{Given a string, count the number of \emph{distinct}
non-empty substrings it contains.}

\begin{patternbox}
\emph{``How many distinct substrings''} is a one-line consequence of the
\textbf{suffix automaton}: each state represents a contiguous range of
substring lengths, so the total distinct count is
$\sum_{v\ne\text{root}}\big(\textit{len}[v]-\textit{len}[\text{link}[v]]\big)$.
\end{patternbox}

\subsection*{Why the length differences count substrings}

Every substring of the string corresponds to exactly one state of the
suffix automaton and one length within that state's range. A state $v$
represents all substrings whose lengths lie in
$(\textit{len}[\text{link}[v]],\ \textit{len}[v]]$ -- that is
$\textit{len}[v]-\textit{len}[\text{link}[v]]$ distinct substrings, each
occurring at $v$'s endpos set. Summing this difference over all non-root
states counts every distinct substring exactly once.

\begin{ideabox}[Each State Owns a Length Interval]
The suffix automaton partitions all distinct substrings among its states,
and within a state the distinct substrings are precisely those of length
strictly greater than the link's length up to the state's own length.
Summing those interval sizes is the distinct-substring count -- no
enumeration required.
\end{ideabox}

\subsection*{Algorithm}

\begin{enumerate}
  \item Build the suffix automaton of the string.
  \item Sum $\textit{len}[v]-\textit{len}[\text{link}[v]]$ over all states
        $v$ except the root.
  \item Output the sum.
\end{enumerate}

\subsection*{Dry run}

\dryrun{$s=$``aa''.}

\begin{itemize}
  \item Distinct substrings: ``a'', ``aa'' -- two of them.
  \item The SAM has states for ``a'' (len $1$, link root len $0$
        $\to1$) and ``aa'' (len $2$, link ``a'' len $1\to1$), summing to
        $2$.
\end{itemize}

\subsection*{C++ implementation}

\begin{lstlisting}
#include <bits/stdc++.h>
using namespace std;

struct SAM {
    struct State { int len, link; map<char,int> next; };
    vector<State> st; int last;
    SAM() { st.push_back({0, -1, {}}); last = 0; }
    void extend(char c) {
        int cur = st.size();
        st.push_back({st[last].len + 1, -1, {}});
        int p = last;
        while (p != -1 && !st[p].next.count(c)) { st[p].next[c] = cur; p = st[p].link; }
        if (p == -1) st[cur].link = 0;
        else {
            int q = st[p].next[c];
            if (st[p].len + 1 == st[q].len) st[cur].link = q;
            else {
                int clone = st.size();
                st.push_back(st[q]);
                st[clone].len = st[p].len + 1;
                while (p != -1 && st[p].next[c] == q) { st[p].next[c] = clone; p = st[p].link; }
                st[q].link = st[cur].link = clone;
            }
        }
        last = cur;
    }
};

int main() {
    string s;
    cin >> s;
    SAM sam;
    for (char c : s) sam.extend(c);

    long long total = 0;
    for (int v = 1; v < (int)sam.st.size(); v++)
        total += sam.st[v].len - sam.st[sam.st[v].link].len;
    cout << total << "\n";
    return 0;
}
\end{lstlisting}

\begin{complexitybox}
\textbf{Time Complexity.} $O(n\log\Sigma)$ to build; the sum is linear in
the number of states. \textbf{Space Complexity.} $O(n)$ states.
\end{complexitybox}

\begin{takeawaybox}
Distinct substrings $=\sum(\textit{len}[v]-\textit{len}[\text{link}[v]])$.
This is the cleanest demonstration that a suffix automaton compresses all
$O(n^2)$ substrings into $O(n)$ states, each owning a length interval. A
suffix array with LCP gives the same count as
$\tfrac{n(n+1)}{2}-\sum\text{LCP}$.
\end{takeawaybox}

\section{Distinct Subsequences}
\label{sec:cses-distinct-subsequences}

\problemstatement{Given a string, count the number of distinct non-empty
subsequences (characters chosen in order, not necessarily contiguous;
duplicates counted once). Output modulo $10^9+7$.}

\begin{patternbox}
\emph{``Count distinct subsequences''} is a linear DP over prefixes. Let
$\textit{dp}[i]$ be the number of distinct subsequences (including the
empty one) of the first $i$ characters; each new character doubles the
count but we must subtract those double-counted by a previous identical
character.
\end{patternbox}

\subsection*{Double, then remove duplicates}

$\textit{dp}[i]$ = distinct subsequences of $s[0..i-1]$, with
$\textit{dp}[0]=1$ (the empty subsequence). Adding character $s[i-1]$, each
existing subsequence can either exclude or include the new character, so
naively $\textit{dp}[i]=2\,\textit{dp}[i-1]$. But if $s[i-1]$ appeared
before at position $j$ (its previous occurrence, 1-indexed), the
subsequences formed by appending this character duplicate exactly those
counted when it last appeared -- namely $\textit{dp}[j-1]$ of them -- so we
subtract:
\[
  \textit{dp}[i]=2\,\textit{dp}[i-1]-\textit{dp}[\,\text{last}[s_i]-1\,].
\]
The final answer is $\textit{dp}[n]-1$ (dropping the empty subsequence).

\begin{ideabox}[Subtract the Last Occurrence's Contribution]
Doubling over-counts precisely the subsequences that already ended in this
character the last time it appeared. Subtracting $\textit{dp}[j-1]$ (the
state just before the previous occurrence) removes exactly those
duplicates -- an inclusion--exclusion in one term.
\end{ideabox}

\subsection*{Algorithm}

\begin{enumerate}
  \item $\textit{dp}[0]=1$; track $\textit{last}[c]$ = previous index of
        character $c$ (0 if none).
  \item For $i=1\ldots n$: $\textit{dp}[i]=2\,\textit{dp}[i-1]$; if
        $s[i-1]$ seen before at $j$, subtract $\textit{dp}[j-1]$; update
        $\textit{last}[s_{i-1}]=i$.
  \item Output $(\textit{dp}[n]-1)\bmod(10^9+7)$.
\end{enumerate}

\subsection*{Dry run}

\dryrun{$s=$``aba''.}

\begin{itemize}
  \item $\textit{dp}[0]=1$; $\textit{dp}[1]=2$ (``'', ``a'').
  \item $\textit{dp}[2]=2\cdot2=4$ (``'',``a'',``b'',``ab''); `b' is new.
  \item $\textit{dp}[3]=2\cdot4-\textit{dp}[0]=8-1=7$ (subtract the
        duplicate from the earlier `a'). Answer $7-1=6$ non-empty
        subsequences.
\end{itemize}

\subsection*{C++ implementation}

\begin{lstlisting}
#include <bits/stdc++.h>
using namespace std;
const long long MOD = 1e9 + 7;

int main() {
    string s;
    cin >> s;
    int n = s.size();
    vector<long long> dp(n + 1, 0);
    dp[0] = 1;
    vector<int> last(26, 0);                       // last[c] = 1-indexed prev pos

    for (int i = 1; i <= n; i++) {
        dp[i] = 2 * dp[i - 1] % MOD;
        int c = s[i - 1] - 'a';
        if (last[c] != 0) dp[i] = (dp[i] - dp[last[c] - 1] + MOD) % MOD;
        last[c] = i;
    }
    cout << (dp[n] - 1 + MOD) % MOD << "\n";        // exclude empty
    return 0;
}
\end{lstlisting}

\begin{complexitybox}
\textbf{Time Complexity.} $O(n)$ -- one pass. \textbf{Space Complexity.}
$O(n)$ for the DP array (reducible to $O(1)$ with rolling values plus the
26 last-occurrence markers).
\end{complexitybox}

\begin{takeawaybox}
Counting distinct subsequences is ``double and subtract the last
occurrence's prefix count.'' The subtraction is a clean inclusion--exclusion
that removes exactly the repeats a duplicate character would create -- a
technique that reappears in distinct-subarray and distinct-path counting.
\end{takeawaybox}

\section{Repeating Substring}
\label{sec:cses-repeating-substring}

\problemstatement{Given a string, find the longest substring that occurs
at least twice (occurrences may overlap). If none exists, report so;
otherwise output one such longest repeating substring.}

\begin{patternbox}
\emph{``Longest substring occurring $\ge2$ times''} falls straight out of
the \textbf{suffix automaton} augmented with endpos sizes: a substring
repeats iff its state's endpos set has size $\ge2$, so the answer is the
longest such state.
\end{patternbox}

\subsection*{Longest state with endpos size at least two}

Recall each SAM state's endpos-set size (computed by propagating counts up
the suffix-link tree) is the number of occurrences of its substrings. A
substring repeats exactly when that count is at least $2$. Among all states
with count $\ge2$, the one with the largest $\textit{len}$ gives the
longest repeating substring; its stored first ending position
(\textit{firstpos}) lets us extract the actual characters.

\begin{ideabox}[Repetition = Endpos Size $\ge 2$]
The suffix automaton already knows how often each substring occurs. So the
longest repeating substring is simply the deepest (largest-\textit{len})
state whose occurrence count is at least two -- no binary search or hashing
needed, and no collision risk.
\end{ideabox}

\subsection*{Algorithm}

\begin{enumerate}
  \item Build the SAM with endpos counts and \textit{firstpos} per state.
  \item Over all states with count $\ge2$, take the one with maximum
        $\textit{len}$.
  \item Output the substring ending at its \textit{firstpos} of that
        length (or state that none exists if no such state).
\end{enumerate}

\subsection*{Dry run}

\dryrun{$s=$``aabaa''.}

\begin{itemize}
  \item ``aa'' occurs at positions ending $2$ and $5$ -- count $2$,
        length $2$.
  \item ``a'' occurs more but is shorter. The longest repeating substring
        is ``aa''.
\end{itemize}

\subsection*{C++ implementation}

\begin{lstlisting}
#include <bits/stdc++.h>
using namespace std;

struct SAM {
    struct State { int len, link, firstpos; map<char,int> next; };
    vector<State> st; vector<long long> cnt; int last;
    SAM() { st.push_back({0, -1, 0, {}}); cnt.push_back(0); last = 0; }
    void extend(char c, int pos) {
        int cur = st.size();
        st.push_back({st[last].len + 1, -1, pos, {}}); cnt.push_back(1);
        int p = last;
        while (p != -1 && !st[p].next.count(c)) { st[p].next[c] = cur; p = st[p].link; }
        if (p == -1) st[cur].link = 0;
        else {
            int q = st[p].next[c];
            if (st[p].len + 1 == st[q].len) st[cur].link = q;
            else {
                int clone = st.size();
                st.push_back(st[q]); cnt.push_back(0);
                st[clone].len = st[p].len + 1;
                while (p != -1 && st[p].next[c] == q) { st[p].next[c] = clone; p = st[p].link; }
                st[q].link = st[cur].link = clone;
            }
        }
        last = cur;
    }
};

int main() {
    string s;
    cin >> s;
    SAM sam;
    for (int i = 0; i < (int)s.size(); i++) sam.extend(s[i], i + 1);

    int m = sam.st.size();
    vector<int> order(m);
    iota(order.begin(), order.end(), 0);
    sort(order.begin(), order.end(),
         [&](int a, int b){ return sam.st[a].len > sam.st[b].len; });
    for (int v : order)
        if (sam.st[v].link >= 0) sam.cnt[sam.st[v].link] += sam.cnt[v];

    int best = 0;
    for (int v = 1; v < m; v++)
        if (sam.cnt[v] >= 2 && sam.st[v].len > sam.st[best].len) best = v;

    if (best == 0) cout << "-1\n";
    else {
        int len = sam.st[best].len, end = sam.st[best].firstpos;
        cout << s.substr(end - len, len) << "\n";
        // (end is 1-indexed end position; substring is s[end-len .. end-1])
    }
    return 0;
}
\end{lstlisting}

\begin{complexitybox}
\textbf{Time Complexity.} $O(n\log\Sigma)$ to build and propagate.
\textbf{Space Complexity.} $O(n)$ states.
\end{complexitybox}

\begin{takeawaybox}
The longest repeating substring is the deepest suffix-automaton state with
occurrence count $\ge2$ -- reusing the endpos-size and first-position
augmentations from earlier sections. A binary-search-plus-hashing solution
also works but risks collisions; the SAM answer is exact.
\end{takeawaybox}

\section{String Functions}
\label{sec:cses-string-functions}

\problemstatement{Given a string, compute two classic string functions:
the \textbf{prefix function} (the border/failure array) and the
\textbf{Z-function}. Output both arrays.}

\begin{patternbox}
This problem is a direct exercise in the two linear-time fundamentals.
The \textbf{prefix function} $\pi[i]$ is the longest proper border of the
prefix ending at $i$; the \textbf{Z-function} $z[i]$ is the longest
substring starting at $i$ that matches a prefix of the whole string.
\end{patternbox}

\subsection*{Two views of the same self-similarity}

Both functions measure how a string matches itself. The prefix function
scans left to right, extending or falling back along the failure chain so
each character is charged $O(1)$ amortised. The Z-function maintains the
rightmost matched interval (the ``Z-box'') and seeds each new index from
its mirror inside that box. They are convertible into each other, but
computing each directly with its own linear routine is simplest.

\begin{ideabox}[Prefix Function: Extend or Fall Back]
At index $i$, try to extend the previous border: if $s[i]$ equals
$s[j]$ (where $j=\pi[i-1]$) increment, else follow $j\gets\pi[j-1]$ until a
match or $j=0$. The Z-function mirror-seeds instead. Both amortise to
$O(n)$ because the fallback/expansion total is bounded by $n$.
\end{ideabox}

\subsection*{Algorithm}

\begin{enumerate}
  \item Compute $\pi$: for each $i$, fall back via $\pi$ until $s[i]$
        matches or the border is empty, then set $\pi[i]$.
  \item Compute $z$: maintain $(l,r)$; seed $z[i]$ from $z[i-l]$ if inside
        the box, then extend by direct comparison.
  \item Output both arrays.
\end{enumerate}

\subsection*{Dry run}

\dryrun{$s=$``abacaba''.}

\begin{itemize}
  \item Prefix function: $[0,0,1,0,1,2,3]$ -- e.g.\ $\pi[6]=3$ because
        ``aba'' is the longest border.
  \item Z-function: $[7,0,1,0,3,0,1]$ -- $z[4]=3$ since ``aba'' at index
        $4$ matches the prefix ``aba''.
\end{itemize}

\subsection*{C++ implementation}

\begin{lstlisting}
#include <bits/stdc++.h>
using namespace std;

vector<int> prefixFunction(const string& s) {
    int n = s.size();
    vector<int> pi(n, 0);
    for (int i = 1; i < n; i++) {
        int j = pi[i - 1];
        while (j > 0 && s[i] != s[j]) j = pi[j - 1];
        if (s[i] == s[j]) j++;
        pi[i] = j;
    }
    return pi;
}

vector<int> zFunction(const string& s) {
    int n = s.size();
    vector<int> z(n, 0);
    z[0] = n;
    int l = 0, r = 0;
    for (int i = 1; i < n; i++) {
        if (i < r) z[i] = min(r - i, z[i - l]);
        while (i + z[i] < n && s[z[i]] == s[i + z[i]]) z[i]++;
        if (i + z[i] > r) { l = i; r = i + z[i]; }
    }
    return z;
}

int main() {
    string s;
    cin >> s;
    for (int x : prefixFunction(s)) cout << x << " ";
    cout << "\n";
    for (int x : zFunction(s)) cout << x << " ";
    cout << "\n";
    return 0;
}
\end{lstlisting}

\begin{complexitybox}
\textbf{Time Complexity.} $O(n)$ for each function. \textbf{Space
Complexity.} $O(n)$ per array.
\end{complexitybox}

\begin{takeawaybox}
The prefix function and Z-function are the two workhorses of string
matching; both run in linear time by never re-comparing a character more
than a constant number of times amortised. Every earlier
border/period/matching section is an application of one of them.
\end{takeawaybox}

\section{Inverse Suffix Array}
\label{sec:cses-inverse-suffix-array}

\problemstatement{Given a suffix array (a permutation $\textit{SA}$ of
$0\ldots n-1$ listing the starting indices of the string's suffixes in
sorted order), reconstruct the \emph{lexicographically smallest} string
over the alphabet whose suffix array is exactly $\textit{SA}$.}

\begin{patternbox}
\emph{``Recover a string from its suffix array''} is a greedy assignment:
walk the suffixes in sorted order and give each position the smallest
character consistent with the ordering, incrementing the character only
when two adjacent sorted suffixes force a strict increase.
\end{patternbox}

\subsection*{When must the character increase?}

Let $\textit{rank}[i]$ be the position of suffix $i$ in $\textit{SA}$.
Process consecutive suffixes $a=\textit{SA}[i-1]$ and $b=\textit{SA}[i]$.
Suffix $a$ must be lexicographically smaller than suffix $b$. Their first
characters satisfy $\text{ch}[a]\le\text{ch}[b]$. We may keep them
\emph{equal} only if the order is already settled by the following
suffixes, i.e.\ $\textit{rank}[a+1]<\textit{rank}[b+1]$ (treating an
out-of-range suffix as the smallest). Otherwise the first characters must
differ, so we increment the current character. Assigning the smallest
letter at $\textit{SA}[0]$ and applying this rule yields the minimal
string.

\begin{ideabox}[Increment Only When Forced]
Keeping characters equal whenever the suffixes after them already
enforce the order gives the smallest possible alphabet usage -- hence the
lexicographically smallest string. The next-suffix rank comparison is the
exact test for ``is a strict increase required here?''
\end{ideabox}

\subsection*{Algorithm}

\begin{enumerate}
  \item Build $\textit{rank}$ from $\textit{SA}$; define
        $\textit{rankNext}(x)=\textit{rank}[x]$ if $x<n$, else $-1$.
  \item Assign character $0$ to $\textit{SA}[0]$. For $i=1\ldots n-1$
        with $a=\textit{SA}[i-1],b=\textit{SA}[i]$: if
        $\textit{rankNext}(a+1)<\textit{rankNext}(b+1)$ keep the current
        character, else increment it; assign it to $b$.
  \item Output the characters (as \texttt{'a'} $+$ value).
\end{enumerate}

\subsection*{Dry run}

\dryrun{$n=4$, $\textit{SA}=[3,1,0,2]$ (suffixes sorted).}

\begin{itemize}
  \item Assign $\text{ch}[3]=a$ (smallest, the sorted-first suffix).
  \item Pair $(3,1)$: $\textit{rankNext}(4)=-1<\textit{rankNext}(2)$, keep
        $\Rightarrow\text{ch}[1]=a$. Pair $(1,0)$: compare
        $\textit{rankNext}(2)$ vs $\textit{rankNext}(1)$; if forced,
        increment.
  \item The construction yields the minimal string consistent with
        $\textit{SA}$.
\end{itemize}

\subsection*{C++ implementation}

\begin{lstlisting}
#include <bits/stdc++.h>
using namespace std;

int main() {
    int n;
    cin >> n;
    vector<int> sa(n);
    for (int& x : sa) cin >> x;              // 0-indexed suffix starts

    vector<int> rankArr(n);
    for (int i = 0; i < n; i++) rankArr[sa[i]] = i;
    auto rankNext = [&](int x) { return x < n ? rankArr[x] : -1; };

    vector<int> ch(n, 0);
    int cur = 0;
    ch[sa[0]] = 0;
    for (int i = 1; i < n; i++) {
        int a = sa[i - 1], b = sa[i];
        if (!(rankNext(a + 1) < rankNext(b + 1))) cur++;   // strict increase forced
        ch[b] = cur;
    }

    string res;
    for (int i = 0; i < n; i++) res += char('a' + ch[i]);
    cout << res << "\n";
    return 0;
}
\end{lstlisting}

\begin{complexitybox}
\textbf{Time Complexity.} $O(n)$ -- one pass building ranks and one pass
assigning characters. \textbf{Space Complexity.} $O(n)$.
\end{complexitybox}

\begin{takeawaybox}
Reconstructing the minimal string from a suffix array is a greedy that
increments a character only when the following suffixes fail to already
order two adjacent entries. This ``smallest label consistent with a given
order'' pattern recurs whenever inverting an ordering into concrete values.
\end{takeawaybox}

\section{String Transform}
\label{sec:cses-string-transform}

\problemstatement{The Burrows--Wheeler transform (BWT) of a string sorts
all rotations of the string (with an appended sentinel \texttt{\#} that is
lexicographically smallest) and outputs the last column. Given a BWT
string, reconstruct the original string.}

\begin{patternbox}
\emph{``Invert the Burrows--Wheeler transform''} is done with the
\textbf{LF-mapping}: from the last column alone you can recover the first
column (it is just the last column sorted), and a rank-based mapping links
each character to its predecessor, letting you walk the original string out
one character at a time.
\end{patternbox}

\subsection*{LF-mapping in one idea}

Let $L$ be the given last column. The first column $F$ is $L$ sorted. The
\textbf{LF-mapping} exploits a key invariant: the $k$-th occurrence of a
character in $L$ corresponds to the $k$-th occurrence of that same
character in $F$. Concretely, for a row $i$, the character $L[i]$ maps to
row $\textit{start}[L[i]]+\textit{rank}_i$, where $\textit{start}[c]$ is
the first row in $F$ beginning with $c$ and $\textit{rank}_i$ is how many
$L[i]$ appeared before index $i$. Following this mapping from the sentinel
row spells the original string in reverse.

\begin{ideabox}[Same Rank in Last and First Columns]
Because equal characters keep their relative order between the sorted-first
and last columns, the LF-mapping is exact: it reconstructs which character
precedes each one in the original. Iterating it $n$ times recovers the
whole string with no ambiguity.
\end{ideabox}

\subsection*{Algorithm}

\begin{enumerate}
  \item Count characters of $L$; compute $\textit{start}[c]$ = number of
        characters smaller than $c$.
  \item Compute $\textit{rank}_i$ for each row and the LF map
        $\textit{lf}[i]=\textit{start}[L[i]]+\textit{rank}_i$.
  \item Start at the sentinel row (row $0$ of $F$), emit $L[\text{row}]$,
        follow $\textit{lf}$; after $n$ steps reverse to get the original.
\end{enumerate}

\subsection*{Dry run}

\dryrun{Conceptually, $L$'s sorted copy is $F$; equal letters keep order.}

\begin{itemize}
  \item From the sentinel row, each LF step jumps to the row whose first
        character is the current $L$ character, at the matching rank.
  \item Emitting characters along this chain and reversing reproduces the
        original string ending in \texttt{\#}.
\end{itemize}

\subsection*{C++ implementation}

\begin{lstlisting}
#include <bits/stdc++.h>
using namespace std;

int main() {
    string L;
    cin >> L;                                    // BWT (last column), contains '#'
    int n = L.size();

    vector<int> cnt(128, 0);
    for (char c : L) cnt[c]++;
    vector<int> start(128, 0);
    for (int c = 1; c < 128; c++) start[c] = start[c - 1] + cnt[c - 1];

    vector<int> lf(n);
    vector<int> occ(128, 0);
    for (int i = 0; i < n; i++) {
        lf[i] = start[(int)L[i]] + occ[(int)L[i]];
        occ[(int)L[i]]++;
    }

    // sentinel '#' is smallest -> first column row 0 begins with '#'
    string res;
    int row = 0;
    for (int i = 0; i < n; i++) {
        res += L[row];
        row = lf[row];
    }
    reverse(res.begin(), res.end());
    cout << res << "\n";                          // original (ending in '#')
    return 0;
}
\end{lstlisting}

\begin{complexitybox}
\textbf{Time Complexity.} $O(n+\Sigma)$ -- linear counting sort plus a
linear LF walk. \textbf{Space Complexity.} $O(n+\Sigma)$.
\end{complexitybox}

\begin{takeawaybox}
Inverting the BWT hinges on the LF-mapping's rank-preservation between the
first and last columns, letting you unwind the original one character at a
time in linear time. The same rank correspondence underlies FM-index
search used in bioinformatics.
\end{takeawaybox}

\section{Substring Order I}
\label{sec:cses-substring-order-i}

\problemstatement{Given a string and an index $k$, find the $k$-th
\emph{distinct} substring in lexicographic order (each distinct substring
counted once).}

\begin{patternbox}
\emph{``$k$-th distinct substring lexicographically''} is a guided walk on
the \textbf{suffix automaton}: precompute, for each state, how many
distinct substrings can be spelled starting from it, then descend edges in
alphabetical order, subtracting counts until $k$ lands on the target.
\end{patternbox}

\subsection*{Counting continuations per state}

Every distinct substring corresponds to a unique path from the initial
state along transitions. Define $d[v]$ = number of distinct substrings
obtainable by continuing from state $v$:
\[
  d[v]=\sum_{(c,u)\ \text{edge from }v}\big(1+d[u]\big),
\]
the $1$ for the single-character extension $c$ and $d[u]$ for everything
that continues past it. Compute $d$ in decreasing order of state length
(edges only go to longer states). To find the $k$-th substring, start at
the root and examine outgoing edges in increasing character order: an edge
to $u$ accounts for $1+d[u]$ substrings. If $k=1$, the current prefix plus
this character is the answer; if $k\le1+d[u]$, append the character, move
to $u$, and decrement $k$ by $1$; otherwise subtract $1+d[u]$ and try the
next character.

\begin{ideabox}[Rank by Skipping Whole Subtrees]
Because $d[u]$ counts \emph{all} distinct substrings through edge $(c,u)$,
we can skip an entire branch by subtracting its total when $k$ is larger --
turning a $k$-th-element search into an $O(\text{answer length}\times
\Sigma)$ descent, never enumerating substrings.
\end{ideabox}

\subsection*{Algorithm}

\begin{enumerate}
  \item Build the SAM; compute $d[v]$ over states in decreasing
        $\textit{len}$.
  \item From the root, repeatedly scan edges in sorted order: locate the
        edge whose cumulative $1+d[u]$ contains $k$; append its character;
        if $k$ becomes $1$ stop, else move into $u$ with $k\gets k-1$.
  \item Output the accumulated string.
\end{enumerate}

\subsection*{Dry run}

\dryrun{$s=$``aba''; distinct substrings sorted: ``a'',``ab'',``aba'',
``b'',``ba''; find $k=4$.}

\begin{itemize}
  \item From root, edge ``a'' covers ``a'',``ab'',``aba'' ($1+d=3$
        substrings). $k=4>3$, subtract $\to k=1$, move to edge ``b''.
  \item Edge ``b'': $k=1\Rightarrow$ answer is ``b''.
\end{itemize}

\subsection*{C++ implementation}

\begin{lstlisting}
#include <bits/stdc++.h>
using namespace std;

struct SAM {
    struct State { int len, link; map<char,int> next; };
    vector<State> st; int last;
    SAM() { st.push_back({0, -1, {}}); last = 0; }
    void extend(char c) {
        int cur = st.size();
        st.push_back({st[last].len + 1, -1, {}});
        int p = last;
        while (p != -1 && !st[p].next.count(c)) { st[p].next[c] = cur; p = st[p].link; }
        if (p == -1) st[cur].link = 0;
        else {
            int q = st[p].next[c];
            if (st[p].len + 1 == st[q].len) st[cur].link = q;
            else {
                int clone = st.size();
                st.push_back(st[q]);
                st[clone].len = st[p].len + 1;
                while (p != -1 && st[p].next[c] == q) { st[p].next[c] = clone; p = st[p].link; }
                st[q].link = st[cur].link = clone;
            }
        }
        last = cur;
    }
};

int main() {
    string s;
    long long k;
    cin >> s >> k;
    SAM sam;
    for (char c : s) sam.extend(c);

    int m = sam.st.size();
    vector<int> order(m);
    iota(order.begin(), order.end(), 0);
    sort(order.begin(), order.end(),
         [&](int a, int b){ return sam.st[a].len > sam.st[b].len; });

    vector<long long> d(m, 0);
    for (int v : order)
        for (auto& [c, u] : sam.st[v].next)
            d[v] += 1 + d[u];                        // distinct continuations

    string res;
    int v = 0;
    while (k > 0) {
        for (auto& [c, u] : sam.st[v].next) {        // sorted char order
            long long here = 1 + d[u];
            if (k <= here) { res += c; k--; v = u; break; }
            k -= here;
        }
    }
    cout << res << "\n";
    return 0;
}
\end{lstlisting}

\begin{complexitybox}
\textbf{Time Complexity.} $O(n\log\Sigma)$ to build and count; the descent
is $O(|\text{answer}|\cdot\Sigma)$. \textbf{Space Complexity.} $O(n)$.
\end{complexitybox}

\begin{takeawaybox}
Ranking distinct substrings is a count-and-descend on the suffix
automaton: $d[v]=\sum(1+d[u])$ lets you skip whole branches. This
``precompute subtree sizes, then walk to the $k$-th'' pattern is the
lexicographic-order workhorse, extended to multiplicity in the next
section.
\end{takeawaybox}

\section{Substring Order II}
\label{sec:cses-substring-order-ii}

\problemstatement{Given a string and an index $k$, find the $k$-th
substring in lexicographic order \emph{counting multiplicity}: a substring
that occurs $t$ times is counted $t$ times, and equal occurrences are
grouped together.}

\begin{patternbox}
This is Substring Order I weighted by occurrences. Each substring's weight
is its \textbf{endpos-set size} (how many times it occurs), so the
per-state continuation count becomes an occurrence-weighted sum, and the
descent subtracts weighted totals.
\end{patternbox}

\subsection*{Weighting by occurrences}

Let $\textit{cnt}[u]$ be the endpos-set size of state $u$ (its occurrence
count, computed by propagating counts up the suffix-link tree). A path
reaching $u$ spells a substring occurring $\textit{cnt}[u]$ times. Define
the occurrence-weighted continuation total:
\[
  f[v]=\sum_{(c,u)\ \text{edge from }v}\big(\textit{cnt}[u]+f[u]\big),
\]
where $\textit{cnt}[u]$ counts the $\textit{cnt}[u]$ occurrences of the
substring ending exactly at $u$, and $f[u]$ counts all longer
continuations with their multiplicities. The descent mirrors Order~I: at an
edge to $u$, if $k\le\textit{cnt}[u]$ the current prefix plus this
character is the answer; else if $k\le\textit{cnt}[u]+f[u]$ move into $u$
after subtracting $\textit{cnt}[u]$; otherwise subtract the whole
$\textit{cnt}[u]+f[u]$ and continue.

\begin{ideabox}[Occurrences Replace the Unit Count]
Order~I used $1$ per distinct substring; Order~II uses $\textit{cnt}[u]$,
the number of occurrences. The identical descent then ranks substrings with
repetition. Computing $\textit{cnt}$ is the same suffix-link propagation
from Counting Patterns (Section~\ref{sec:cses-counting-patterns}).
\end{ideabox}

\subsection*{Algorithm}

\begin{enumerate}
  \item Build the SAM; compute $\textit{cnt}[u]$ (endpos sizes) via
        link-tree propagation.
  \item Compute $f[v]=\sum_{(c,u)}(\textit{cnt}[u]+f[u])$ in decreasing
        $\textit{len}$ order.
  \item Descend from the root: at edge to $u$, if $k\le\textit{cnt}[u]$
        output and stop; else if $k\le\textit{cnt}[u]+f[u]$ append, move to
        $u$, subtract $\textit{cnt}[u]$; else subtract $\textit{cnt}[u]+
        f[u]$.
\end{enumerate}

\subsection*{Dry run}

\dryrun{$s=$``aa''; substrings with multiplicity, sorted:
``a''(pos 1), ``a''(pos 2), ``aa''; find $k=2$.}

\begin{itemize}
  \item $\textit{cnt}$ of state ``a'' is $2$; of ``aa'' is $1$.
  \item Edge ``a'' to state ``a'': $\textit{cnt}=2\ge k=2\Rightarrow$
        answer ``a'' (its second occurrence).
\end{itemize}

\subsection*{C++ implementation}

\begin{lstlisting}
#include <bits/stdc++.h>
using namespace std;

struct SAM {
    struct State { int len, link; map<char,int> next; };
    vector<State> st; vector<long long> cnt; int last;
    SAM() { st.push_back({0, -1, {}}); cnt.push_back(0); last = 0; }
    void extend(char c) {
        int cur = st.size();
        st.push_back({st[last].len + 1, -1, {}}); cnt.push_back(1);
        int p = last;
        while (p != -1 && !st[p].next.count(c)) { st[p].next[c] = cur; p = st[p].link; }
        if (p == -1) st[cur].link = 0;
        else {
            int q = st[p].next[c];
            if (st[p].len + 1 == st[q].len) st[cur].link = q;
            else {
                int clone = st.size();
                st.push_back(st[q]); cnt.push_back(0);
                st[clone].len = st[p].len + 1;
                while (p != -1 && st[p].next[c] == q) { st[p].next[c] = clone; p = st[p].link; }
                st[q].link = st[cur].link = clone;
            }
        }
        last = cur;
    }
};

int main() {
    string s;
    long long k;
    cin >> s >> k;
    SAM sam;
    for (char c : s) sam.extend(c);

    int m = sam.st.size();
    vector<int> order(m);
    iota(order.begin(), order.end(), 0);
    sort(order.begin(), order.end(),
         [&](int a, int b){ return sam.st[a].len > sam.st[b].len; });
    for (int v : order)
        if (sam.st[v].link >= 0) sam.cnt[sam.st[v].link] += sam.cnt[v];

    vector<long long> f(m, 0);
    for (int v : order)
        for (auto& [c, u] : sam.st[v].next)
            f[v] += sam.cnt[u] + f[u];               // occurrence-weighted

    string res;
    int v = 0;
    while (k > 0) {
        for (auto& [c, u] : sam.st[v].next) {
            long long occ = sam.cnt[u];
            if (k <= occ) { res += c; k = 0; break; }          // land here
            if (k <= occ + f[u]) { res += c; k -= occ; v = u; break; }
            k -= occ + f[u];
        }
    }
    cout << res << "\n";
    return 0;
}
\end{lstlisting}

\begin{complexitybox}
\textbf{Time Complexity.} $O(n\log\Sigma)$ build and counts; descent
$O(|\text{answer}|\cdot\Sigma)$. \textbf{Space Complexity.} $O(n)$.
\end{complexitybox}

\begin{takeawaybox}
Order~II is Order~I with the unit count replaced by each state's occurrence
count. The two problems together show how one weighting choice on the same
count-and-descend framework toggles between distinct and multiplicity
ranking.
\end{takeawaybox}

\section{Substring Distribution}
\label{sec:cses-substring-distribution}

\problemstatement{Given a string of length $n$, for every length
$L=1,2,\ldots,n$ output the number of \emph{distinct} substrings of that
length.}

\begin{patternbox}
\emph{``Distinct substrings of each length''} uses the same
\textbf{suffix automaton} fact as distinct-substring counting, but resolved
per length: each state contributes exactly one distinct substring to every
length in its interval
$(\textit{len}[\text{link}],\ \textit{len}]$, so a difference array over
those intervals gives the per-length counts.
\end{patternbox}

\subsection*{Each state covers a contiguous length range}

A suffix automaton state $v$ represents distinct substrings of every length
from $\textit{len}[\text{link}[v]]+1$ up to $\textit{len}[v]$, one distinct
substring per length (they are the longer substring's suffixes, all sharing
$v$'s endpos class). Therefore the number of distinct substrings of length
$L$ equals the number of states whose interval contains $L$. Add $+1$ over
each state's interval with a difference array, then take a prefix sum to
read off the count for every length.

\begin{ideabox}[Difference Array Over State Intervals]
Instead of iterating each length inside each state ($O(n^2)$), mark
$\textit{diff}[\textit{len}[\text{link}]+1]\mathrel{+}{=}1$ and
$\textit{diff}[\textit{len}+1]\mathrel{-}{=}1$ per state, then prefix-sum
once. Every state's whole length interval is applied in $O(1)$.
\end{ideabox}

\subsection*{Algorithm}

\begin{enumerate}
  \item Build the suffix automaton.
  \item For each non-root state $v$: $\textit{diff}[\textit{len}[
        \text{link}[v]]+1]\mathrel{+}{=}1$;
        $\textit{diff}[\textit{len}[v]+1]\mathrel{-}{=}1$.
  \item Prefix-sum $\textit{diff}$; output entries for $L=1\ldots n$.
\end{enumerate}

\subsection*{Dry run}

\dryrun{$s=$``aba''.}

\begin{itemize}
  \item Distinct substrings by length: length $1$: ``a'',``b'' $\to2$;
        length $2$: ``ab'',``ba'' $\to2$; length $3$: ``aba'' $\to1$.
  \item The interval contributions sum, via the difference array, to
        $2,2,1$.
\end{itemize}

\subsection*{C++ implementation}

\begin{lstlisting}
#include <bits/stdc++.h>
using namespace std;

struct SAM {
    struct State { int len, link; map<char,int> next; };
    vector<State> st; int last;
    SAM() { st.push_back({0, -1, {}}); last = 0; }
    void extend(char c) {
        int cur = st.size();
        st.push_back({st[last].len + 1, -1, {}});
        int p = last;
        while (p != -1 && !st[p].next.count(c)) { st[p].next[c] = cur; p = st[p].link; }
        if (p == -1) st[cur].link = 0;
        else {
            int q = st[p].next[c];
            if (st[p].len + 1 == st[q].len) st[cur].link = q;
            else {
                int clone = st.size();
                st.push_back(st[q]);
                st[clone].len = st[p].len + 1;
                while (p != -1 && st[p].next[c] == q) { st[p].next[c] = clone; p = st[p].link; }
                st[q].link = st[cur].link = clone;
            }
        }
        last = cur;
    }
};

int main() {
    string s;
    cin >> s;
    int n = s.size();
    SAM sam;
    for (char c : s) sam.extend(c);

    vector<long long> diff(n + 2, 0);
    for (int v = 1; v < (int)sam.st.size(); v++) {
        int lo = sam.st[sam.st[v].link].len + 1;    // shortest length in v
        int hi = sam.st[v].len;                     // longest length in v
        diff[lo] += 1;
        diff[hi + 1] -= 1;
    }
    long long run = 0;
    for (int L = 1; L <= n; L++) { run += diff[L]; cout << run << " "; }
    cout << "\n";
    return 0;
}
\end{lstlisting}

\begin{complexitybox}
\textbf{Time Complexity.} $O(n\log\Sigma)$ to build; $O(n)$ for the
difference array and output. \textbf{Space Complexity.} $O(n)$.
\end{complexitybox}

\begin{takeawaybox}
Each suffix-automaton state owns a contiguous interval of substring
lengths; a difference array turns ``count distinct substrings per length''
into one linear sweep. This closes the chapter's suffix-automaton arc:
presence, count, position, distinct count, ranking, and now the full
length distribution -- all from one structure.
\end{takeawaybox}

